%!TEX root = ../dissertation.tex
\begin{savequote}[75mm]
The only true wisdom is in knowing you know nothing.
\qauthor{Socrates}
\end{savequote}

\chapter{Discussion}
\label{ch:discussion}

This thesis studies whether the geometric structure of prompt embeddings carries information about hallucination. The answer that emerges is not simply that geometry predicts hallucination, but that hallucination itself has a \emph{difficulty structure}. Some hallucinations are easy to fix: a well-chosen system prompt eliminates them. Others resist every intervention we tested. That difficulty is not random; it is systematic, and it can be read from the geometry of the question before the model generates a single token. This is the central finding of the thesis, and it reframes hallucination from a binary failure mode (the model got it wrong) to a graded property of the input space (some questions are harder to get right, and we can tell which ones in advance).

The geometric evidence for this claim operates at two scales. At the category level, the separation is unsurprising: category membership---whether a question concerns a nonexistent entity, a factual recall, or a logical impossibility---accounts for most of the predictive power (cross-validated AUC of 0.78). But the benchmark was designed with a difficulty gradient across categories, so this separation is by construction. The non-trivial finding is within-category. Holding question type constant, embedding density still distinguishes prompts that cause hallucination from those answered correctly ($p < 10^{-5}$, both models, Bonferroni-surviving, cross-validated AUC $\approx 0.67$ for the nonexistent category). The same density feature predicts which hallucinations resist prompting ($p = 0.046$/$0.012$) and fine-tuning ($p = 0.0017$, $|r| = 0.80$). 

The practical arc of the thesis follows from this structure. All four system-prompt interventions significantly reduce hallucination, and the result replicates at five times the original benchmark scale. But the more consequential finding is that prompt responses themselves become raw material for permanent model improvement. The best response per prompt across all conditions becomes the fine-tuning training target, producing training data that is over 99\% correct. This reframes prompt engineering from a runtime band-aid to a data generation strategy---a way to produce the training signal that makes models permanently safer.

Fine-tuning on that curated data consolidates the cautious behavior into model weights: the fine-tuned models match the best prompt prefix without any runtime system prompt, and hallucination drops by nearly 90\%. The cost is increased refusals on factual questions. But when the model does regress on prompts it previously handled correctly, the vast majority of those regressions are refusals rather than new fabrications. In other words, the model errs toward caution, not toward new kinds of error. That caution is domain-specific: on TruthfulQA, an external benchmark that tests a different type of hallucination entirely, our fine-tuned model shows significant improvement with no increase in refusal. Three further tests---invariance to template diversity, transfer to held-out category types, and minimal inflation from entity-level train-test overlap---confirm that the model learns a general behavioral strategy rather than memorizing patterns from its training data. The full arc is a single integrated pipeline: geometry predicts risk, prompts generate a corrective training signal, geometry identifies which prompts resist correction, and fine-tuning consolidates the rest. For practitioners, the takeaway is that models can be made permanently safer through curated self-distillation, with the cost both predictable and domain-specific.

This thesis deliberately prioritizes depth over breadth: a controlled benchmark, two well-characterized models, and one embedding space, examined thoroughly with category-level controls and multiple robustness checks. That depth comes at the cost of generality. All geometric findings are correlational. Density predicts hallucination and intervention resistance with consistent direction across models and captures more signal than surface-level text properties. A confound such as entity rarity independently driving both sparse embeddings and poor model recall remains possible. The intervention pipeline covers two open-weight models and one embedding space; the cross-model benchmark spans 10 models, but prompting and fine-tuning have not been tested on dense architectures, larger models, or closed-source systems. A preliminary comparison across three embedding models (Appendix~\ref{app:embedding-robustness}) finds that centrality is stable, but density's correlation with hallucination varies. Whether the within-category density signal transfers to other embeddings is the most important open question for the geometric claims in this thesis. All prompts are template-generated, providing experimental control at the cost of ecological validity; whether findings hold on free-form user queries remains open, and we note that all experiments are conducted in English. 

Several of these findings connect to established frameworks in cognitive science. We do not claim that LLMs implement the same mechanisms as human cognition, but the structural parallels are worth noting. Humans confabulate (or produce plausible but false memories) when probed at the boundaries of their knowledge, and LLMs can exhibit the same structure. Hallucinations are essentially plausible fabrication concentrated at knowledge boundaries rather than in regions of complete knowledge or complete ignorance. The parallel is particularly concrete for our plausible-fake category, which mirrors the Deese-Roediger-McDermott paradigm in false memory research \citep{roediger1995creating}: participants confidently ``remember'' semantically plausible lures that were never presented, just as models confidently describe fabricated entities designed to sound real. More broadly, the thesis asks whether models can be made to ``know what they don't know'' \citep{kadavath2022language}---a question that cognitive science studies under the heading of metacognition. Our geometric features serve as an external metacognitive signal; fine-tuning teaches the model to act on that signal by expressing uncertainty when its knowledge is unreliable. The precision-recall tradeoff can be understood through signal detection theory \citep{green1966signal}: the fine-tuned model shifts its criterion toward caution, reducing false alarms at the cost of increased misses, with the shift targeted to entity-existence judgments rather than applied globally. These parallels are interpretive rather than empirical, but the structural similarities may reflect general properties of systems that must act under incomplete knowledge.

Two directions for future work seem especially promising. The first is causal. Density predicts where hallucinations occur and where they resist correction; the open question is whether it plays a causal role. If prompts were rephrased to shift their embeddings into denser neighborhoods, would hallucination rates decrease? Such a perturbation experiment would distinguish density as a genuine factor from a correlate of entity rarity, and would open the door to geometry-based intervention. The second direction concerns the relationship between external and internal representations. This thesis finds that an external embedding model's density predicts hallucination across two architecturally distinct target models, implying the signal lives in the question text itself. But do the models' own internal representations carry a stronger version of the same signal? If model-internal density is a better predictor, the external density we measured is a lower bound on what is geometrically available. If external density works equally well, hallucination risk is genuinely a property of the input space, independent of how any particular model encodes it. Either answer would sharpen our understanding of what density captures and how far the geometric diagnostic extends.

The broader motivation for this work, beyond hallucination specifically, is that the geometry of representational spaces may carry general information about system reliability. The principle that emerges here is that sparse neighborhoods predict unreliable behavior, and that this unreliability has structure: some failures are fixable, others are not. This is not inherently specific to language models. Any system that must act on learned representations faces regions where those representations are thin, and understanding where they thin out is a prerequisite for knowing when to trust the system's outputs. The specific tools developed in this thesis---within-category geometric controls, best-per-prompt curation, density as a pre-generation difficulty signal---may find application wherever input geometry and output reliability intersect.

This thesis asks whether the structure of a model's representational space reveals where its knowledge thins out. That structure predicts not only which prompts cause hallucination, but which hallucinations resist correction. And the behavioral caution learned from this signal generalizes across entities, templates, and question types. Much remains to be understood, but the core finding is clear: hallucination has a difficulty structure, and that structure is legible from the geometry of the input space.
